% GENERATED FILE. Edit canonical lessons, then run npm run generate:books.
% canonical-source-hash: fnv1a64:eccb7489a6beba78
% canonical-lessons: SA-C299-ghosayati, SA-C299-sankatam, SA-C299-savadhanah, SA-C299-nivarayati, SA-C299-apattih

\chapter{A Warning}
\label{ch:sa-a-warning}
\hlchaptermodality{299}

\begin{chapteropening}
\textbf{By the end of this chapter:} \emph{I can pass on a warning and say what to be careful of.}

Complete the last lesson of chapter 299: I can pass on a warning and say what to be careful of.
\end{chapteropening}

\section[\texorpdfstring{\sk{घोषयति} (ghoṣayati)}{घोषयति (ghoṣayati)}]{\sk{घोषयति} (ghoṣayati) — announces}

\label{lesson:SA-C299-ghosayati}



\begin{quote}
Before the new one: say the Sanskrit for a citizen, then the Sanskrit for what is left.
\end{quote}

\subsection*{You'll want to know: \sk{घोषयति}}
\textbf{\sk{घोषयति}} — \emph{ghoṣayati} — \textquotedblleft{}announces\textquotedblright{}.

\begin{quote}
\sk{अधिकारी} \sk{उच्चैः} \sk{घोषयति।}
\end{quote}

\begin{quote}
\emph{adhikārī uccaiḥ ghoṣayati.}
\end{quote}

\begin{quote}
The official announces it loudly.
\end{quote}

Past (he …): \sk{सः} \sk{उच्चैः} \sk{घोषितवान्}, \emph{saḥ uccaiḥ ghoṣitavān}.

\begin{cousinweb}
\textbf{\sk{घोषयति}} (\emph{ghoṣayati}) comes from the root √ghuṣ, \textquotedblleft{}to sound, cry aloud, announce publicly\textquotedblright{}; \emph{ghoṣayati} is its causative (MW).
\end{cousinweb}

\begin{grammarlens}[title={What you've built}]
One more word for telling what happened.
\end{grammarlens}

\subsection*{Your turn}
\begin{itemize}
\raggedright
  \item \emph{Say it:} \emph{ghoṣayati}
  \item \emph{Say it:} \emph{ghoṣayati}, once more
  \item \emph{Say it:} \emph{adhikārī uccaiḥ ghoṣayati.} — The official announces it loudly.
  \item \emph{Recall:} say the Sanskrit for a citizen, then the Sanskrit for what is left, then say \emph{ghoṣayati} again
\end{itemize}

\begin{tcolorbox}[breakable,colback=teal!4,colframe=teal!35!black,title={Before you move on}]
What does \sk{घोषयति} mean? (\textquotedblleft{}Announces\textquotedblright{}.) Say it once more.
\end{tcolorbox}
\section[\texorpdfstring{\sk{संकटम्} (saṅkaṭam)}{संकटम् (saṅkaṭam)}]{\sk{संकटम्} (saṅkaṭam) — danger}

\label{lesson:SA-C299-sankatam}



\begin{quote}
Before the new one: say the Sanskrit for what is left, then the Sanskrit for announces.
\end{quote}

\subsection*{You'll want to know: \sk{संकटम्}}
\textbf{\sk{संकटम्}} — \emph{saṅkaṭam} — \textquotedblleft{}danger\textquotedblright{}.

\begin{quote}
\sk{नद्याः} \sk{समीपे} \sk{संकटम्} \sk{अस्ति।}
\end{quote}

\begin{quote}
\emph{nadyāḥ samīpe saṅkaṭam asti.}
\end{quote}

\begin{quote}
There is danger near the river.
\end{quote}

A neuter noun.

\begin{cousinweb}
\textbf{\sk{संकटम्}} (\emph{saṅkaṭam}) MW thinks it is probably a Prakrit form of \emph{saṃkṛta}, \textquotedblleft{}brought together\textquotedblright{}, built from \emph{sam} and the root √kṛ; its senses run from narrow to dangerous (MW).
\end{cousinweb}

\begin{grammarlens}[title={What you've built}]
One more word for telling what happened.
\end{grammarlens}

\subsection*{Your turn}
\begin{itemize}
\raggedright
  \item \emph{Say it:} \emph{saṅkaṭam}
  \item \emph{Say it:} \emph{saṅkaṭam}, once more
  \item \emph{Say it:} \emph{nadyāḥ samīpe saṅkaṭam asti.} — There is danger near the river.
  \item \emph{Recall:} say the Sanskrit for what is left, then the Sanskrit for announces, then say \emph{saṅkaṭam} again
\end{itemize}

\begin{tcolorbox}[breakable,colback=teal!4,colframe=teal!35!black,title={Before you move on}]
What does \sk{संकटम्} mean? (\textquotedblleft{}Danger\textquotedblright{}.) Say it once more.
\end{tcolorbox}
\section[\texorpdfstring{\sk{सावधानः} (sāvadhānaḥ)}{सावधानः (sāvadhānaḥ)}]{\sk{सावधानः} (sāvadhānaḥ) — careful}

\label{lesson:SA-C299-savadhanah}



\begin{quote}
Before the new one: say the Sanskrit for announces, then the Sanskrit for danger.
\end{quote}

\subsection*{You'll want to know: \sk{सावधानः}}
\textbf{\sk{सावधानः}} — \emph{sāvadhānaḥ} — \textquotedblleft{}careful\textquotedblright{}.

\begin{quote}
\sk{भवान्} \sk{सावधानः} \sk{भवतु।}
\end{quote}

\begin{quote}
\emph{bhavān sāvadhānaḥ bhavatu.}
\end{quote}

\begin{quote}
Please be careful.
\end{quote}

\begin{cousinweb}
\textbf{\sk{सावधानः}} (\emph{sāvadhānaḥ}) is built from \emph{sa}, \textquotedblleft{}with\textquotedblright{}, and \emph{avadhāna}, \textquotedblleft{}attention\textquotedblright{}; MW gives \emph{sāvadhāna} as attentive, heedful, careful (MW).
\end{cousinweb}

\begin{grammarlens}[title={What you've built}]
One more word for telling what happened.
\end{grammarlens}

\subsection*{Your turn}
\begin{itemize}
\raggedright
  \item \emph{Say it:} \emph{sāvadhānaḥ}
  \item \emph{Say it:} \emph{sāvadhānaḥ}, once more
  \item \emph{Say it:} \emph{bhavān sāvadhānaḥ bhavatu.} — Please be careful.
  \item \emph{Recall:} say the Sanskrit for announces, then the Sanskrit for danger, then say \emph{sāvadhānaḥ} again
\end{itemize}

\begin{tcolorbox}[breakable,colback=teal!4,colframe=teal!35!black,title={Before you move on}]
What does \sk{सावधानः} mean? (\textquotedblleft{}Careful\textquotedblright{}.) Say it once more.
\end{tcolorbox}
\section[\texorpdfstring{\sk{निवारयति} (nivārayati)}{निवारयति (nivārayati)}]{\sk{निवारयति} (nivārayati) — keeps away; prevents}

\label{lesson:SA-C299-nivarayati}



\begin{quote}
Before the new one: say the Sanskrit for danger, then the Sanskrit for careful.
\end{quote}

\subsection*{You'll want to know: \sk{निवारयति}}
\textbf{\sk{निवारयति}} — \emph{nivārayati} — \textquotedblleft{}keeps away; prevents\textquotedblright{}.

\begin{quote}
\sk{रक्षकः} \sk{जनान्} \sk{नद्याः} \sk{निवारयति।}
\end{quote}

\begin{quote}
\emph{rakṣakaḥ janān nadyāḥ nivārayati.}
\end{quote}

\begin{quote}
The guard keeps people away from the river.
\end{quote}

Past (he …): \sk{सः} \sk{जनान्} \sk{निवारितवान्}, \emph{saḥ janān nivāritavān}.

\begin{cousinweb}
\textbf{\sk{निवारयति}} (\emph{nivārayati}) is built from \emph{ni} and the causative of the root √vṛ, \textquotedblleft{}to ward off\textquotedblright{}; MW gives \emph{ni-vārayati} as to hold back from, prevent and stop (MW).
\end{cousinweb}

\begin{grammarlens}[title={What you've built}]
One more word for telling what happened.
\end{grammarlens}

\subsection*{Your turn}
\begin{itemize}
\raggedright
  \item \emph{Say it:} \emph{nivārayati}
  \item \emph{Say it:} \emph{nivārayati}, once more
  \item \emph{Say it:} \emph{rakṣakaḥ janān nadyāḥ nivārayati.} — The guard keeps people away from the river.
  \item \emph{Recall:} say the Sanskrit for danger, then the Sanskrit for careful, then say \emph{nivārayati} again
\end{itemize}

\begin{tcolorbox}[breakable,colback=teal!4,colframe=teal!35!black,title={Before you move on}]
What does \sk{निवारयति} mean? (\textquotedblleft{}Keeps away; prevents\textquotedblright{}.) Say it once more.
\end{tcolorbox}
\section[\texorpdfstring{\sk{आपत्तिः} (āpattiḥ)}{आपत्तिः (āpattiḥ)}]{\sk{आपत्तिः} (āpattiḥ) — a misfortune}

\label{lesson:SA-C299-apattih}



\begin{quote}
Before the new one: say the Sanskrit for careful, then the Sanskrit for keeps away.
\end{quote}

\subsection*{You'll want to know: \sk{आपत्तिः}}
\textbf{\sk{आपत्तिः}} — \emph{āpattiḥ} — \textquotedblleft{}a misfortune\textquotedblright{}.

\begin{quote}
\sk{सहसा} \sk{ग्रामे} \sk{आपत्तिः} \sk{आगच्छति।}
\end{quote}

\begin{quote}
\emph{sahasā grāme āpattiḥ āgacchati.}
\end{quote}

\begin{quote}
Suddenly misfortune comes to the village.
\end{quote}

A feminine noun.

\begin{cousinweb}
\textbf{\sk{आपत्तिः}} (\emph{āpattiḥ}) is built from \emph{ā} and the root √pad, \textquotedblleft{}to fall\textquotedblright{}: MW gives \emph{āpatti} as happening, incurring and misfortune or calamity (MW).
\end{cousinweb}

\begin{grammarlens}[title={What you've built}]
One more word for telling what happened.
\end{grammarlens}

\subsection*{Your turn}
\begin{itemize}
\raggedright
  \item \emph{Say it:} \emph{āpattiḥ}
  \item \emph{Say it:} \emph{āpattiḥ}, once more
  \item \emph{Say it:} \emph{sahasā grāme āpattiḥ āgacchati.} — Suddenly misfortune comes to the village.
  \item \emph{Recall:} say the Sanskrit for careful, then the Sanskrit for keeps away, then say \emph{āpattiḥ} again
\end{itemize}

\begin{tcolorbox}[breakable,colback=teal!4,colframe=teal!35!black,title={Before you move on}]
What does \sk{आपत्तिः} mean? (\textquotedblleft{}A misfortune\textquotedblright{}.) Say it once more.
\end{tcolorbox}
